\subsection{(\S5) The wizard of Oz: a vision of development in the American political psyche}

\begin{enumerate}

\item \ovalbox{\it ``The Wizard of Oz'' is a 1939 American musical.} The Wizard
of Oz, the 1939 MGM color musical starring Judy Garland, Bert Lahr, and Ray
Bolger (Wizard of Oz, 1939), has been reissued so often that its presence
within the American psyche has the quality of a recurrent dream.

\item \ovalbox{\it how to examine this musical from Jungian perspective}
Dorothy's individuation proceeds through her links to the other characters.
The interaction of eight characters is one way to visualize the differentiation
of consciousness in an individuating psyche.

\item \ovalbox{\it section title explained} In a film with particular resonance
such as The Wizard of Oz, these roles can be specifically metaphoric of the
political developments emerging in the American psyche at the moment of the
film's construction.

\item \ovalbox{\it ?} Indeed, part of the comedy of the Lion, as played by Bert
Lahr, is that his intuition is so inferior that, for all his tearfulness, he
never seems to anticipate the real danger he invites by the physical
(extraverted sensation) abandon with which he wades into each new situation. It
is not so much that he lacks courage as that his courage is of the puer
aeternus type. Bert Lahr’s Lion displays an inflated sense of what he can pull
off that is easily punctured by reality, at which point he falls into despair
over his own cowardice.

\item \ovalbox{\it ?} Thomson’s simile echoes the sensibility of The Wizard of
Oz taken as a whole. The consciousness of the characters and their relative
positions combine to give the film itself a definite consciousness—that is, a
certain set of concerns that define it not only as a psychological, but also a
political space. By establishing the positions of the characters in relation to
each other in terms of the way they feel about themselves, how the other
characters relate to them, and the degree of emotional development they
display, we can understand the consciousness politics of the film.

\item \ovalbox{\it ?} Like Dorothy, we are being initiated into the politics of
confrontation, where consciousnesses encounter opposing attitudes, with respect
to introversion and extraversion, in the form of figures they experience as
shadowy. The characters that give the film its spine of integrity can be
diagrammed in relation to their antagonists, naming the psychological types
associated with each character. We should understand that even a figure of the
unconscious can exhibit a characteristic type, implying that there is
intelligence in the unconscious, although it is typically used in an
unconscious and frequently unethical way, hence the description of such
functions of potential but problematic consciousness as ‘in shadow’ (Figure
5.7). In this scheme, we can see that Dorothy’s heroic extraverted feeling is
opposed by the introverted feeling of Almira Gulch, and that the Scarecrow’s
righteous indignation, expressing an animus demand for a higher standard of
integrity, is directed against the pompous obfuscations of the Wizard.

\end{enumerate}
